%! Setting Up the SVD — Unit 7, Lesson 7.0 — Guided Notes
\documentclass[10pt]{article}
\usepackage{linalg-article}
\usepackage{linalg-boxes}
\newcommand{\TallMath}[1]{$\displaystyle #1\rule[-1.4em]{0pt}{3.2em}$}
\newcommand{\uu}{\mathbf{u}}
\newcommand{\vv}{\mathbf{v}}
\newcommand{\ee}{\mathbf{e}}
\newcommand{\zero}{\mathbf{0}}
\newcommand{\T}{^{\top}}

\begin{document}

\pageheader{Unit 7, Lesson 7.0}{Guided Notes}
\namedateperiod

\vspace{0.12in}

\begin{objectivebox}
\textbf{By the end of this lesson, I will be able to\dots}
\begin{itemize}\setlength{\itemsep}{2pt}
  \item explain that a matrix maps the \emph{unit circle} to an \emph{ellipse}, whose semi-axis lengths are the \emph{singular values};
  \item find a perpendicular input pair $\vv_1,\vv_2$ sent to a perpendicular output pair with $A\vv_i=\sigma_i\uu_i$, and read the singular values $\sigma_i$;
  \item state the SVD $A=U\Sigma V\T$ and say why every matrix has one.
\end{itemize}
\end{objectivebox}

\vspace{0.08in}

\begin{vocabbox}
\small
Fill in each term as we build it together.
\vspace{2pt}
\termblanklong{Singular value $\sigma$ (a stretch factor)}
\termblanklong{Singular vectors $\vv_i$ (input) and $\uu_i$ (output)}
\termblanklong{The relation $A\vv_i=\sigma_i\uu_i$}
\termblanklong{The SVD $A=U\Sigma V\T$}
\end{vocabbox}

\begin{hookbox}
In Unit~5 a matrix \emph{moved} the whole plane; in Unit~6 we found the few directions it only
\emph{stretched}. Now watch what a matrix does to the \textbf{unit circle} (all vectors of length $1$).
\begin{itemize}\setlength{\itemsep}{1pt}
  \item What shape do you think the circle becomes after you multiply every point by a matrix? \writeline
  \item If the circle becomes a stretched oval, what two numbers would describe how far it stretches? \writeline
\end{itemize}
\end{hookbox}

\begin{notesbox}{1. A Matrix Turns the Unit Circle into an Ellipse}
Multiply every vector of length $1$ (the unit circle) by a matrix and the circle becomes an
\textbf{ellipse}. The simplest case is a \emph{diagonal} matrix, which just stretches the axes. Take
$D=\begin{bsmallmatrix} 3 & 0\\ 0 & 2 \end{bsmallmatrix}$:
\[
  D\begin{bmatrix} 1\\ 0 \end{bmatrix}=\begin{bmatrix}\rule{0pt}{1.1em}\blank{0.8cm}\\[2pt]\blank{0.8cm}\end{bmatrix}
  \ (\text{length }\blank{0.8cm}),\qquad
  D\begin{bmatrix} 0\\ 1 \end{bmatrix}=\begin{bmatrix}\rule{0pt}{1.1em}\blank{0.8cm}\\[2pt]\blank{0.8cm}\end{bmatrix}
  \ (\text{length }\blank{0.8cm}).
\]
So the circle becomes an ellipse stretched to a half-width of \blank{0.8cm} along the $x$-axis and
\blank{0.8cm} along the $y$-axis. Those two half-widths --- the \textbf{singular values} --- are
$\sigma_1=\blank{0.8cm}$ and $\sigma_2=\blank{0.8cm}$ (we list the larger one first).
\end{notesbox}

\begin{notesbox}{2. The Big Idea: Perpendicular In $\to$ Perpendicular Out}
Most matrices are not diagonal, so they \emph{tilt} the ellipse. The key fact of Unit~7:
\begin{center}
\emph{every matrix sends some perpendicular pair of input directions to a perpendicular pair of
output directions} --- only stretched.
\end{center}
In symbols, there are unit vectors $\vv_1\perp\vv_2$ and $\uu_1\perp\uu_2$ with
\[
  A\vv_1=\sigma_1\uu_1,\qquad A\vv_2=\sigma_2\uu_2,\qquad \sigma_1\ge\sigma_2\ge0.
\]
The $\vv$'s are the \textbf{input (right) singular vectors}, the $\uu$'s the \textbf{output (left)
singular vectors}, and the $\sigma$'s the \textbf{singular values}. Take the Unit~6 matrix
$A=\begin{bsmallmatrix} 2 & 1\\ 1 & 2 \end{bsmallmatrix}$ with its two perpendicular directions
$\vv_1=\begin{bsmallmatrix} 1\\ 1 \end{bsmallmatrix}$, $\vv_2=\begin{bsmallmatrix} 1\\ -1 \end{bsmallmatrix}$:
\[
  A\vv_1=\begin{bmatrix}\rule{0pt}{1.1em}\blank{0.8cm}\\[2pt]\blank{0.8cm}\end{bmatrix}
  =\blank{0.8cm}\cdot\vv_1,\qquad
  A\vv_2=\begin{bmatrix}\rule{0pt}{1.1em}\blank{0.8cm}\\[2pt]\blank{0.8cm}\end{bmatrix}
  =\blank{0.8cm}\cdot\vv_2.
\]
Perpendicular in: $\vv_1\cdot\vv_2=\blank{0.8cm}$. \; Perpendicular out:
$(A\vv_1)\cdot(A\vv_2)=\blank{0.8cm}$. \; So the stretch factors --- the singular values --- are
$\sigma_1=\blank{0.8cm}$ and $\sigma_2=\blank{0.8cm}$. (As \emph{unit} vectors, divide each direction
by $\sqrt2$; here the output directions equal the input directions.)

\begin{center}
\begin{tikzpicture}[scale=0.5]
  % LEFT: unit circle with perpendicular v's
  \begin{scope}[shift={(-4.4,0)}]
    \draw[->, charcoal, thin] (-1.8,0)--(1.8,0);
    \draw[->, charcoal, thin] (0,-1.8)--(0,1.8);
    \draw[burgundy!70, thick] (0,0) circle (1);
    \draw[->, ultra thick, burgundy] (0,0)--(0.707,0.707) node[above right=-2pt]{\scriptsize $\vv_1$};
    \draw[->, ultra thick, royalblue] (0,0)--(0.707,-0.707) node[below right=-2pt]{\scriptsize $\vv_2$};
    \node[charcoal, below] at (0,-2.0) {\scriptsize unit circle};
  \end{scope}
  % ARROW
  \draw[->, very thick, goldacc] (-2.0,0)--(-0.2,0) node[midway, above]{\scriptsize $A$};
  % RIGHT: tilted ellipse with perpendicular u's
  \begin{scope}[shift={(3.1,0)}]
    \draw[->, charcoal, thin] (-3.4,0)--(3.4,0);
    \draw[->, charcoal, thin] (0,-3.4)--(0,3.4);
    \draw[burgundy!70, thick, rotate=45] (0,0) ellipse (3 and 1);
    \draw[->, ultra thick, burgundy] (0,0)--(2.121,2.121) node[above right=-2pt]{\scriptsize $A\vv_1=3\uu_1$};
    \draw[->, ultra thick, royalblue] (0,0)--(0.707,-0.707) node[below right=-2pt]{\scriptsize $A\vv_2=\uu_2$};
    \node[charcoal, below] at (0,-3.6) {\scriptsize ellipse: semi-axes $\sigma_1=3,\ \sigma_2=1$};
  \end{scope}
\end{tikzpicture}
\end{center}
\end{notesbox}

\begin{notesbox}{3. Reading It Off --- Input and Output Frames Can Differ}
For a symmetric matrix (above) the input and output directions happened to match. In general they are
\emph{different} perpendicular frames. Take the swap-and-stretch
$S=\begin{bsmallmatrix} 0 & 2\\ 1 & 0 \end{bsmallmatrix}$ acting on the axes $\ee_1,\ee_2$:
\[
  S\begin{bmatrix} 1\\ 0 \end{bmatrix}=\begin{bmatrix}\rule{0pt}{1.1em}\blank{0.8cm}\\[2pt]\blank{0.8cm}\end{bmatrix}
  \ (\text{length }\blank{0.8cm}),\qquad
  S\begin{bmatrix} 0\\ 1 \end{bmatrix}=\begin{bmatrix}\rule{0pt}{1.1em}\blank{0.8cm}\\[2pt]\blank{0.8cm}\end{bmatrix}
  \ (\text{length }\blank{0.8cm}).
\]
The inputs $\ee_1,\ee_2$ are perpendicular; the outputs are perpendicular too
($\begin{bsmallmatrix} 0\\1 \end{bsmallmatrix}\cdot\begin{bsmallmatrix} 2\\0 \end{bsmallmatrix}=\blank{0.8cm}$).
Listing the bigger stretch first, the singular values are $\sigma_1=\blank{0.8cm}$ and
$\sigma_2=\blank{0.8cm}$. Notice the largest input direction $\vv_1=\begin{bsmallmatrix} 0\\1 \end{bsmallmatrix}$
is sent to the \emph{different} output direction $\uu_1=\begin{bsmallmatrix} 1\\0 \end{bsmallmatrix}$ ---
input and output frames need not be the same.
\end{notesbox}

\begin{notesbox}{4. The SVD $A=U\Sigma V\T$ --- and Why It Matters}
Collect the pieces: put the input unit vectors $\vv_1,\vv_2$ into the columns of $V$, the output unit
vectors $\uu_1,\uu_2$ into the columns of $U$, and the singular values down the diagonal of $\Sigma$.
Then every matrix factors as
\[
  A = U\,\Sigma\,V\T \qquad(\Sigma=\begin{bsmallmatrix} \sigma_1 & 0\\ 0 & \sigma_2 \end{bsmallmatrix}),
\]
the \textbf{Singular Value Decomposition}: rotate to the input frame ($V\T$), stretch by the
$\sigma$'s ($\Sigma$), rotate to the output frame ($U$).

\vspace{2pt}
\textbf{Why it matters.} Singular values are always \blank{2.6cm} (they are lengths), and
\emph{every} matrix has an SVD --- even non-square ones --- while eigenvectors need a square matrix and
may not exist. That is why the SVD is the tool behind image compression and data science.

\vspace{2pt}
\textbf{How do we \emph{find} them} without guessing? The singular values are $\sigma=\sqrt{\lambda}$ for
the eigenvalues $\lambda$ of $A\T A$ --- a symmetric matrix (Unit~6). That is Lesson~7.1.

\vspace{2pt}
\textbf{The road ahead.} \textbf{7.1} \emph{finding} singular values/vectors from $A\T A$\; $\bullet$\;
\textbf{7.2} compressing images\; $\bullet$\; \textbf{7.3} principal component analysis\; $\bullet$\;
\textbf{7.4} the victory of orthogonality.
\end{notesbox}

\begin{practicebox}
Show your work.
\begin{enumerate}[leftmargin=1.6em, itemsep=4pt]
  \item For $D=\begin{bsmallmatrix} 5 & 0\\ 0 & 2 \end{bsmallmatrix}$, what ellipse does the unit circle become? Give the singular values $\sigma_1,\sigma_2$. \writeline
  \item For $S=\begin{bsmallmatrix} 0 & 3\\ 1 & 0 \end{bsmallmatrix}$, compute $S\ee_1$ and $S\ee_2$, check the outputs are perpendicular, and read $\sigma_1,\sigma_2$. \writeline
  \item Why can a singular value never be negative? \writeline
\end{enumerate}
\end{practicebox}

\end{document}
